
\section{Space and Orthogonality}

As established, the Cognitive Subject, supported by the computational capacity
of sedimented Past, possesses the faculty to transform diachronic sequences
into synchronic presence. Sequential facts $s \in \mathbb{F}$ are
observed `at once' within a single Synthetic Interstice.

To derive Spatial extension from this pure information, we must return to the
primitive operation of \textbf{Composition} (\texttt{AND}, $\oplus$) based on
co-existence.

\begin{formalism}[Spatial Point Discernible]

    A \textbf{Spatial Point Discernible} $\delta_p \in \mathbb{E}$ is a
    composed of numeric discernibles. 
    
    \begin{equation}
        \delta_p := n_1 \oplus n_2 \oplus \dots \oplus n_d 
    \end{equation}
    
    The cardinality $d$ of the fused numeric discernibles intrinsically
    establishes the \textbf{Dimensionality} of the point. Because these
    discernibles co-exist in the Now without the sequential operator ($\succ$),
    they exhibit zero logical distance between them.

\end{formalism}

The discernible of a point $\delta_p = n_1 \oplus n_2 \dots \oplus n_d$
neither possesses temporal expansion nor spatial expansion while being
associated with $\top$, i.e.\ the event $e_p$. Once fixated (associated with
$\bot$), though, it is equipped with temporal stability as a trace in the
sediment of tangible Past, but not with spatial expansion.

\subsection{Space and the Absence of Redundancy}

We now consider the relation of two distinct composite discernibles within the
purview of the Subject. The faculty of Differentiation implies an active event
of difference between the two. Since the Synthetic Interstice structurally
prohibits redundancy, this implied event must naturally carry the information
about their non-equivalence in its absolute purest, redundancy-free form. 

Let $\ominus$ denote the logical operation of \textit{Differentiation} between
discernibles (the complement to Composition $\oplus$). Consider a differentiating
discernible concerning two composite discernibles $\delta_a$ and $\delta_b$,
given by:

\begin{equation}
\begin{aligned}
    \delta_a &= \delta_{a,1} \oplus \delta_{a,2} \\
    \delta_b &= \delta_{b,1} \oplus \delta_{b,2} 
\end{aligned}
\end{equation}

The difference between them cannot ontologically exist as a redundant, uncompressed
composite discernible:

\begin{equation}
    \delta_{c,redundant} = (\delta_{a,1} \ominus \delta_{b,1}) \oplus (\delta_{a,2} \ominus \delta_{b,2}) 
\end{equation}

if it can be mathematically expressed in a strictly purer state:

\begin{equation}
    \delta_{c} = \delta'_a \ominus \delta'_b 
\end{equation}

where $(\delta'_a \ominus \delta'_b)$ carries the total magnitude of the difference
between $\delta_a$ and $\delta_b$ along a single informational axis.

\begin{formalism}[Stretch Discernible]

    A \textbf{Stretch Discernible} $t \in \mathbb{F}$ is the conveys
    a redundancy-free differentiation between two composite discernibles
    $\delta_a$ and $\delta_b$ belonging to two events $e_a$ and $e_b$. It
    maintains a single discernible $(\delta'_a \ominus \delta'_b)$ representing
    the total modality of their difference. 

    \begin{equation}
        t := \langle \bot, (\delta'_a \ominus \delta'_b) \oplus C \rangle \in \mathbb{F}
    \end{equation}

    Here, $C$ represents the \textbf{Shared Core}---the dimensions of the
    composition that have been rendered equivalent and mutually redundant between
    the two points.

\end{formalism}

Let us focus on specific kinds of discernibles, namely `Numeric Discernibles'.
The difference between two arbitrary composite numeric discernibles---that is,
$d$-dimensional vectors $\delta_a=(\delta_{a,1}, \dots, \delta_{a,d})$ and
$\delta_b=(\delta_{b,1}, \dots, \delta_{b,d})$---can be exhaustively described
by the difference of a single pair of transformed discernibles $\delta'_a$ and
$\delta'_b$. This establishes the ontological state of a `Stretch'. 

This Stretch is the explicit expression of a redundancy collapse in the
differentiation between two vectors. 

Mathematically, this redundancy collapse operates via the formal principles of
\textit{Spectral Decomposition} or \textit{Principal Axis
Transformation}~\cite{Euler}. It is established by `translating' the composite
discernible into an equivalent informational basis such that the entirety of
the differentiation is expressed by a single magnitude $\mathbb{D}$ in the form
$(\mathbb{D}, 0, \dots, 0)$---a perfect projection of the difference onto a
single principal axis. The remaining discernibles (the zeros) are rendered
structurally redundant and collapse instantly into the shared core
$C$.

It is of profound metaphysical significance that this exact
mathematical operation constitutes the transcendental genesis of spatial
extension. This ontological necessity provides a formal grounding for why
Eigenvalue analysis and Spectral Decomposition exhibit such ``unreasonable
effectiveness''~\cite{Wigner1960} and ubiquitous utility across the empirical
sciences. They are not merely descriptions of an external physical world; they
are the foundational informational algorithms by which the Cognitive Subject
compresses and structures perceivable reality.

The absolute Absence of Redundancy therefore necessitates this transformational
logic as a fundamental structural faculty of the Synthetic Interstice.

\subsection{The Line Assertion}

A Stretch Discernible, maintains the difference of two composite discernibles
in a single discernible $\mathbb{D}$ while maintaining a redundant shared core $C$.
With $C$ as a virtual seed and $\mathbb{D}$ as a element of identity, the conditions
for a \textbf{Generative Set} are met. 

While the direct experience of an infinite number of points with the same core
$C$ is structurally impossible within the domain of conceptions, the labeling
of a significant magnitude of recurrence as `Universal' may occur (again). That
is, an assertion may be made that a generative set exists consisting of an 
infinite number of points with a shared core $C$. 

\begin{formalism}[Line Assertion]

    A \textbf{Line Assertion} is the assertion $l \in \mathbb{A}^\diamond$
    that, behind a metric fact $m \in \mathbb{A}^\top$, that repetition
    expands to infinity, i.e.\ that there is an infinite number of points
    that share the same $C$ considering expansion and detail. 

    That includes the assertion that there would exist an entity in the Noumena
    that governs the regularity.

    Let $t = \langle \bot, \mathbb{D} \oplus C \rangle$ be a Stretch Fact and $m$ 
    be the Metric Fact representing the finite recurrence of this equivalent 
    difference. The Line Assertion $l \in \mathbb{A}^\diamond$ is the speculative 
    fixation defined as:

    \begin{equation}
        l := \langle \bot, (\varnothing, \infty(\mathbb{D}) \oplus C) \rangle \in \mathbb{A}^\diamond
    \end{equation}

\end{formalism}

As with the Causal Assertion, the assertion of an entity in the
un-conceptualizable Noumena, exposes $\varnothing$, wrongness. We identified
the origin of wrongness in the cognitive subject as Will. Therefore, the
postulation of continuous space is a profound act of the \textbf{Will} ($\Pi$).
The Cognitive Subject constructs the spatial continuum to structure and
navigate the infinitely complex discernibles of the Absolute. As with the Causal
Assertion, the un-conceptualizable nature of the Noumena prevents the
perceiption of a conceptional answer to why the Noumena does not instantiate 
events that immediately refute the Line Assertion.

